\chapter{Introduction}

%  It’s a good idea to try to write the introduction to your final report early on in the project. However, you will find it hard, as you won’t yet have a complete story and you won’t know what your main contributions are going to be. However, the exercise is useful as it will tell you what you don’t yet know and thus what questions your project should aim to answer. For the interim report this section should be a short, succinct, summary of the project’s main objectives. Some of this material may be re-usable in your final report, but the chances are that your final introduction will be quite different.  You are therefore advised to keep this part of the interim report short, focusing on the following questions:

% What is the problem,
% why is it interesting,
% what’s your main idea for solving it?
% (DON'T use those three questions as subheadings however!  The answers should emerge from what you write.)

Testing software is important. Across the U.S. economy, the cost of poor software quality was estimated at at least \$2.41 trillion in 2022, with technical debt around \$1.52 trillion - illustrating the economic risk of defects and the value of rigorous verification and validation \cite{globenewswire_2022_cisq,synopsysinc_2022_software}. However, maintaining a high-quality test suite is challenging: it is difficult to enumerate all relevant inputs, states, and corner cases, and infeasible to exercise every potential path through complex software systems. Systematic techniques like fuzzing help by exploring vast input spaces - including malformed or unexpected inputs that traditional test design may miss - and can be run continuously with runtime feedback (e.g., sanitizers, coverage) to catch regressions early \cite{mans_2021_the,yu_2024_fuzzing}.

Fuzzing has proven to be a very successful technique to find bugs and vulnerabilities in code. Some of the most important software we use - language compilers, operating system components, cryptography libraries, and database systems - have had critical bugs discovered and fixed via fuzzing \cite{mans_2021_the}. For example, Google’s large-scale fuzzing infrastructure has reported thousands of security vulnerabilities and tens of thousands of functional bugs across hundreds of open-source projects, leveraging coverage-guided engines (e.g., libFuzzer, AFL++) with sanitizers and scalable orchestration \cite{googlesecurityblog_2023_aipowered}. These systems rely on code-coverage guidance to steer execution toward new edges and behaviours; while coverage cannot literally guarantee “all code paths” are explored, it provides a practical, measurable heuristic to maximize path diversity during a campaign \cite{googlesecurityblog_2023_aipowered,googlesecurityblog_2024_leveling}.

There are two major categories of fuzzers: generation-based and mutation-based. Generation-based fuzzers use input formats or grammars to synthesize inputs from scratch that satisfy syntactic rules of the SUT. Mutation-based fuzzers begin with a corpus of example inputs (seeds) and iteratively mutate them \cite{mans_2021_the}. In practice, grey-box mutation-based fuzzers often dominate modern vulnerability discovery because they (a) start from real inputs that already respect many syntactic/semantic constraints, enabling rapid progress past shallow checks, and (b) use lightweight program feedback (e.g., edge coverage, value profiling, branch distances) to adapt mutations toward novel paths and deeper states - especially effective for uncovering unexpected or undefined behaviour in complex code bases \cite{mans_2021_the,wang_2020_a}.

Despite the strengths of mutation-based fuzzers, their effectiveness often hinges on the quality and variety of seed inputs. A narrow or unrepresentative corpus can trap a campaign in local minima, limiting path diversity and bug discovery; conversely, a rich, diverse seed set boosts pass rates through parsers and unlocks deeper program states \cite{yu_2024_fuzzing}. Manually crafting such corpora is labour‑intensive and time‑consuming, and using existing samples “from the wild” offers no guarantee of exercising the full behavioural envelope of the SUT. As a result, many fuzzers invest in seed curation, deduplication, minimization, and dynamic corpus management to keep exploration efficient and targeted \cite{wang_2020_a}.

With the advent and widespread adoption of large language models (LLMs), their reasoning and contextual knowledge can alleviate several setbacks of existing fuzzing approaches. In mainstream software assurance, early evidence from large-scale industrial deployments shows that AI assistance can expand test generation and accelerate the on-boarding of new targets to continuous fuzzing; for example, Google reports that AI-generated fuzz targets and helpers have led to more vulnerabilities being found, and expanded coverage across diverse ecosystems \cite{googlesecurityblog_2024_leveling,googlesecurityblog_2024_scaling}. In fuzzing research, LLMs have been used to generate grammars and structured inputs, to mutate inputs with semantic awareness, and to broaden language and domain coverage - spanning compilers, APIs, protocol/file formats, and more \cite{xia_2023_universal,yang_2023_whitebox,chen_2025_elfuzz,wang_2020_a}.

Work has been done in using LLMs to generate fuzzers themselves, guided by coverage data \cite{chen_2025_elfuzz}, as well as generating inputs directly with a simple crash oracle \cite{xia_2023_universal,lin_2025_lmfuzz}. However, we have not seen any projects using LLMs to generate inputs, but also in a coverage guided manner. Our hypothesis is that using LLMs to generate and evolve seed inputs directly in a coverage guided manner will result in a more semantically and structurally diverse corpus, which will in turn lead to better fuzzing performance.

We present \toolname, a novel fuzzing approach that leverages LLMs to generate and evolve seed inputs directly in a coverage guided manner. \toolname uses an LLM to synthesize initial seed inputs that are semantically rich and structurally varied, and then iteratively mutates these seeds using the LLM, guided by coverage feedback from the SUT. This allows \toolname to explore deeper program states and uncover more complex bugs than traditional mutation-based fuzzers that rely on simpler mutation strategies.
